%%%%%%%%%%%%%%%%%%%%%%%%%%%%% Define Article %%%%%%%%%%%%%%%%%%%%%%%%%%%%%%%%%%
\documentclass{article}
%%%%%%%%%%%%%%%%%%%%%%%%%%%%%%%%%%%%%%%%%%%%%%%%%%%%%%%%%%%%%%%%%%%%%%%%%%%%%%%

%%%%%%%%%%%%%%%%%%%%%%%%%%%%% Using Packages %%%%%%%%%%%%%%%%%%%%%%%%%%%%%%%%%%
\usepackage{listings, xcolor}
\usepackage{parskip}
\usepackage{hyperref}
\usepackage{amsmath}
\usepackage{graphicx}
\usepackage[a4paper, margin = 1.5cm]{geometry}
\usepackage{float}
\usepackage{amssymb}

%%%%%%%%%%%%%%%%%%%%%%%%%%%%%%% Title & Author %%%%%%%%%%%%%%%%%%%%%%%%%%%%%%%%
\title{Parte 8/9}
\author{Trentin Tommaso}
%%%%%%%%%%%%%%%%%%%%%%%%%%%%%%%%%%%%%%%%%%%%%%%%%%%%%%%%%%%%%%%%%%%%%%%%%%%%%%%

\begin{document}
  \section{Modelli delle reti}
    Modelli generativi: processi stocastici che generano / producono grafi. Se i grafi generati sono simili alle reti reali, possiamo:
    \begin{enumerate}
      \item Ipotizzare che stessi processi generativi siano alla base delle reti reali;
      \item Studiare le reti reali in modo più efficiente / efficace;
      \item Classificare le reti reali in varie tipologie (che consentono poi di giustificare / spiegare fenomeni). 
    \end{enumerate}
  \section{Reti Regolari}
    Reti regolari = reti con topologia regolare, costruita a tavolino (non casuale e non stocastica).
    \subsection{Proprietà}
      \begin{itemize}
        \item \textbf{Distanze} (average path length e diametro): spesso \textbf{grandi };
        \item Coefficiente di \textbf{clustetring}: spesso \textbf{alto};
        \item \textbf{Distribuzione dei gradi}: spesso con un \textbf{picco} (e spesso 0 altrove);
        \item \textbf{Connetività}: spesso \textbf{completamente connesse}.
      \end{itemize}
  \section{Reti Casuali}
    Grafo con \textit{n} nodi e archi "a caso", individualmente (ogni arco creato indipendentemente dagli altri). Assunzione più semplice possibile, corrisponde ad assumere che amicizie (legami un una rete sociale) nascano a caso.
    \subsection{Modello G(n,m)}
      \begin{itemize}
        \item Grafo ha \textit{n nodi} (fissati); 
        \item Grafo ha \textit{m archi} piazzati a caso fra gli n nodi;
        \item Grafo solitamente senza multi-archi o auto-archi;
        \item Grafo solitamente indiretto. 
      \end{itemize}
      Si studiano il caso $n\rightarrow\infty$:"$Lim_(n\rightarrow\infty)G(n,m)$", "Limite termodinamico" (grafi grandi) ed il \textbf{caso medio} che succede in media su un insieme di tanti ($\infty$) grafi casuali. Si studiano in quanto comodi analiticamente (almeno per $n\rightarrow\infty$) e in quanto permettono di trovare proprietà \textbf{tipiche} (più interessanti, non casi particolari), solitamente più rappresentative. 
    \subsection{Modello G(n,p)}
      Alcune proprietà di G(n,m) sono facilmente calcolabili (es Grado medio $\langle K\rangle = 2m / n$), altre no\dots Si usa quindi di più l'altro modello \textbf{G(n,p)}:
      \begin{itemize}
        \item Più "maneggiabile" matematicamente e analiticamente;
        \item n nodi;
        \item Considero indipendentemente ogni arco "potenziale"; 
        \item L'arco esiste con una \textbf{probabilità p}.
      \end{itemize}
      In pratica, al posto di fissare il # di archi, si fissa la probabilità (indipendente) di esistenza archi, il numero effettivo di archi sarà in media: $$m=p*\frac{n(n-1)}{2}=\binom{n}{2}p$$ Il grado medio invece viene calcolato con: $$\lang K\rang=\frac{2\lang m\rang}{n}=(n-1)p$$ solitamente denotato con $$c=(n-1)p$$
      \subsection{}
\end{document}
